% !PRINTED 375
\weberpage{375}
%<<CONT
determination of an eighth double tangent from a given complete system. But we do better to determine three simultaneously. So let now

\weq{1}{x_1,\ x_2,\ x_3,\ x_4,\ x_5,\ x_6,\ x_7}

be a complete seven-system assumed as known. We imagine the three \gesperrt{Steiner} complexes formed, in which six at a time of the double tangents (1) are contained, with the exclusion first of $x_1$, then of $x_2$, lastly of $x_3$.

The 15 double tangents newly entering there we denote by the letter $\xi$, and obtain these three complexes [by \S.\,99, (1), (2)] in the form:

\weq{2}{\begin{aligned}
x_2\xi_3 &,\ x_3\xi_2,\ x_4\xi_{41},\ x_5\xi_{51},\ x_6\xi_{61},\ x_7\xi_{71}\\
x_3\xi_1 &,\ x_1\xi_3,\ x_4\xi_{42},\ x_5\xi_{52},\ x_6\xi_{62},\ x_7\xi_{72}\\
x_1\xi_2 &,\ x_2\xi_1,\ x_4\xi_{43},\ x_5\xi_{53},\ x_6\xi_{63},\ x_7\xi_{73},
\end{aligned}}

and by the notation of \S.\,99 we have

\[
\xi_1 = [2\,3],\ \xi_2 = [3\,1],\ \xi_3 = [1\,2],\ \xi_{41} = [1\,4],\ \ldots
\]

From (2) there results yet another \gesperrt{Steiner} complex, which forms a syzygetic pair with each of the complexes (2):

\weq{3}{x_1\xi_1,\ x_2\xi_2,\ x_3\xi_3,\ \ldots}

and this complex contains none of the double tangents $x_4, x_5, x_6, x_7$.

If now we take the functions sought, $\xi_1, \xi_2, \xi_3$, multiplied by the suitable constant factors, we can put the equation of the curve of the fourth order, by \S.\,96, (9), into the form

\weq{4}{\sqrt{x_1\xi_1} + \sqrt{x_2\xi_2} + \sqrt{x_3\xi_3} = 0}

and thus take the rational function $f$, which, set equal to zero, gives the equation of the curve, in each of the three mutually identical forms

\weq{5}{f = 4x_2\xi_2x_3\xi_3 - u_1^2 = 4x_3\xi_3x_1\xi_1 - u_2^2 = 4x_1\xi_1x_2\xi_2 - u_3^2,}

where

\weq{6}{\begin{aligned}
u_1 &= -x_1\xi_1 + x_2\xi_2 + x_3\xi_3\\
u_2 &= x_1\xi_1 - x_2\xi_2 + x_3\xi_3\\
u_3 &= x_1\xi_1 + x_2\xi_2 - x_3\xi_3
\end{aligned}}

is set. Now [by (2)] $x_2\xi_3\, x_4\xi_{41}$ also form a syzygetic quadruple, and consequently we can, if over a
%CONT>>
